\chapter{调试、测试与可靠性工程}
\label{chap:debug-reliability}

嵌入式故障经常跨越电源、时钟、信号、处理器、内核和应用边界。
有效调试依赖可观测性、可复现条件和逐层缩小范围，而不是反复修改代码等待现象消失。
可靠性工程则要求在交付前主动构造极端负载、资源耗尽、掉电和器件异常，证明系统能够检测、隔离并恢复故障。

\section{分层故障模型}

建议按以下层次建立假设：
\begin{enumerate}
  \item \textbf{供电与复位}：电压、纹波、时序、欠压、复位脉宽和启动模式；
  \item \textbf{时钟与信号}：晶振、PLL、总线时序、电平、终端和信号完整性；
  \item \textbf{处理器执行}：取指地址、异常向量、栈、MMU、缓存和内存；
  \item \textbf{系统软件}：Bootloader、内核、驱动、文件系统和服务依赖；
  \item \textbf{应用与数据}：协议状态、资源泄漏、并发、配置和持久化一致性。
\end{enumerate}

每次实验应只改变一个关键变量，并保留镜像摘要、配置、硬件版本、环境条件和完整日志。

\section{硬件观测工具}

\subsection{万用表与电子负载}

万用表适合确认稳态电压、电阻、短路和连续性，不能替代示波器观察瞬态跌落和高频纹波。
电子负载可验证电源限流、负载瞬态和热稳定性。测量低电阻或大电流路径时，应考虑表笔和走线压降，并优先使用四线法或专用电流采样。

\subsection{示波器}

示波器测量的可信度取决于探头带宽、接地回路、输入阻抗、采样率和触发条件。
观察电源纹波时应使用短接地弹簧和适当带宽限制；观察高速边沿时，长鳄鱼夹接地线本身可能产生足以误导判断的振铃。

常见触发对象包括欠压、异常复位、总线错误、过长中断脉冲和启动时序。应同时观测因果相关信号，例如电源轨、复位、时钟和故障 GPIO。

\subsection{逻辑分析仪与协议分析}

逻辑分析仪适合捕获数字时序和长时间协议事件，但其阈值模型不能反映模拟电平质量。
协议解码结果必须结合原始波形、采样率和触发位置审查，避免把边沿失真或电气冲突误判为纯协议错误。

\section{从复位到用户空间的调试路径}

\subsection{无输出系统}

当系统没有串口输出时，推荐依次确认：
\begin{enumerate}
  \item 电源轨和复位是否满足数据手册要求；
  \item 启动 strap、存储器和 Boot ROM 下载模式是否正确；
  \item JTAG/SWD 能否连接，处理器停在何处；
  \item 异常向量、栈指针和程序计数器是否位于有效区域；
  \item 串口的时钟、引脚复用、电平和波特率是否匹配；
  \item DRAM 初始化前后执行路径是否发生变化。
\end{enumerate}

不要在尚未证明 DRAM 稳定时依赖复杂日志和动态内存。早期启动应保留最小串口、状态 GPIO 或片上 SRAM 日志路径。

\subsection{Linux 启动问题}

Linux 启动调试应保留 Bootloader 环境、内核命令行、设备树摘要和从最早 console 开始的日志。
发生挂起时，可通过 early console、initcall 调试、动态调试或临时 tracepoint 确定阶段。
如果最小 initramfs 能启动，而正式根文件系统不能启动，问题通常已收敛到存储、挂载、用户空间依赖或安全策略。

QEMU 可以快速验证内核配置、启动参数和部分驱动逻辑，但模拟器不能证明真实电源、时钟、DMA、一致性和外设时序正确。

\subsection{QEMU 与 GDB 内核调试}

可重复的内核调试环境至少固定 QEMU machine、CPU 型号、内存容量、内核配置、DTB、rootfs 和命令行。内核应保留 \texttt{CONFIG\_DEBUG\_INFO}，并用未剥离的 \texttt{vmlinux} 提供符号；QEMU 实际加载的压缩镜像与 GDB 使用的符号文件必须来自同一次构建。

\begin{lstlisting}[language=bash,caption={QEMU 停机等待 GDB 的启动方式（AArch64 示例）}]
qemu-system-aarch64 \
  -machine virt -cpu cortex-a57 -m 1024 -nographic \
  -kernel Image -initrd rootfs.cpio \
  -append 'console=ttyAMA0 earlycon nokaslr' \
  -S -gdb tcp::1234
\end{lstlisting}

\begin{lstlisting}[language=bash,caption={GDB 连接内核}]
gdb-multiarch vmlinux
(gdb) target remote :1234
(gdb) hbreak start_kernel
(gdb) continue
\end{lstlisting}

关闭 KASLR 便于初始学习和符号定位，不能据此形成发布配置。调试模块时还需要根据目标系统中的实际加载地址执行 \texttt{add-symbol-file}；调试早期页表切换时，应区分链接虚拟地址、加载物理地址和当前 MMU 状态。

QEMU 适合验证启动流程、异常路径、锁与部分虚拟设备驱动。涉及缓存一致性、DMA 时序、中断丢失、PHY、模拟电气和真实闪存故障时，必须回到目标硬件和硬件在环环境复验。

\section{软件调试与追踪}

\subsection{固件调试}

GDB、JTAG/SWD 和硬件断点适合控制执行与检查寄存器。观察实时系统时要注意停核可能改变外设、看门狗和多核时序。
对于难以复现的故障，应使用环形事件记录、硬件 trace 或在复位后保留的故障区，而不是只依赖交互式单步。

异常处理程序至少应保存异常原因、程序计数器、栈指针、关键寄存器、任务标识和固件版本。
解析工具必须使用与目标镜像完全对应的 ELF 和映射文件。

\subsection{Linux 故障信息}

内核 oops 和 panic 的价值取决于符号、配置和镜像匹配。
产品应归档未剥离的 \texttt{vmlinux}、模块符号、System.map、配置和源码提交。
pstore/ramoops 可以跨重启保存有限的崩溃信息；kdump 适用于资源允许且需要完整内存转储的系统。

性能与时序问题可使用：
\begin{itemize}
  \item \textbf{ftrace/trace-cmd}：函数、调度、中断和内核事件；
  \item \textbf{perf}：硬件计数器、采样和调用栈；
  \item \textbf{eBPF}：受控地观测内核和用户空间事件；
  \item \textbf{lockdep、KASAN、KCSAN、UBSAN}：锁、内存和并发错误检测；
  \item \textbf{kmemleak 与专用统计}：泄漏和资源生命周期分析。
\end{itemize}

诊断配置会显著改变内存、性能和时序。发现问题后必须在接近发布的配置中复验。

\section{测试层次}

\subsection{静态检查与单元测试}

编译告警、静态分析、格式检查、设备树 schema、配置检查和链接布局检查应作为提交门禁。
纯算法、协议解析和状态机应尽量从硬件访问中分离，使其能够在主机环境进行快速单元测试和模糊测试。

\subsection{集成测试与硬件在环}

集成测试验证真实驱动、服务和外设之间的契约。硬件在环（Hardware-in-the-Loop, HIL）系统可以控制电源、复位、输入信号、网络和测量设备，并自动收集串口、功耗和测试结果。

HIL 工装必须自身可校准、可版本化，并能区分被测设备故障和测试基础设施故障。

\subsection{压力与耐久测试}

压力测试应覆盖比标称负载更恶劣但仍在设计边界内的条件：
\begin{itemize}
  \item CPU、内存、存储和网络并发满载；
  \item 高频中断、队列拥塞和连接反复建立；
  \item 文件系统接近满载及日志持续写入；
  \item 长时间挂起恢复、热插拔和服务重启；
  \item 温度、电压和时钟容差边界；
  \item 多日或多周运行中的泄漏、计数器溢出和时间漂移。
\end{itemize}

测试结果应记录最大值和分布，而不仅是平均值。实时延迟和启动时间通常更关心最坏值或高分位数。

\section{故障注入}

故障注入用于验证错误路径，而不是破坏性地随机操作系统。每个注入点应对应一个预期检测机制、超时、告警和恢复结果。

常见场景包括：
\begin{itemize}
  \item I2C/SPI/CAN 超时、CRC 错误和设备无响应；
  \item DMA 未完成、中断丢失或中断风暴；
  \item 内存分配失败、队列满和文件描述符耗尽；
  \item 存储读写失败、坏块增加和文件系统损坏；
  \item 网络断开、DNS 失败、证书过期和服务器返回异常；
  \item 在镜像写入、元数据提交和首次启动期间断电；
  \item 关键任务卡死、进程崩溃和看门狗复位。
\end{itemize}

故障恢复必须具有时间上限。无限重试既可能阻塞系统，也会放大器件故障和网络拥塞。

\section{看门狗与健康管理}

看门狗只能证明负责喂狗的执行路径仍在运行，不能自动证明系统健康。
可靠设计通常由健康管理任务汇总关键任务心跳、调度延迟、资源水位和外设状态，只有所有必要条件满足时才刷新硬件看门狗。

看门狗超时必须大于经过验证的最长正常阻塞时间，同时小于产品允许的故障恢复时间。
复位后应持久化复位原因和有限诊断上下文，并限制反复重启：连续失败达到阈值时进入恢复或安全模式。

\section{可靠性设计}

\subsection{失效安全与降级}

系统应定义故障时的安全状态。例如停止执行器、保持最后安全输出、关闭高功率通道或切换备用传感器。
“继续运行”不是默认正确策略；应根据危害分析决定停机、降级或冗余接管。

\subsection{资源与寿命预算}

可靠性预算包括 RAM/Flash 裕量、CPU 峰值、存储写放大、擦写寿命、日志速率、电源余量和热裕量。
监控指标应与这些预算直接对应，并设置趋势告警，避免只有资源完全耗尽后才暴露问题。

\section{发布验收清单}

\begin{itemize}
  \item 是否能够根据版本号找到对应源码、配置、符号、设备树和镜像摘要？
  \item 是否定义并测量了启动时间、实时延迟、资源峰值和功耗上限？
  \item 是否完成关键接口、满载、长稳、低功耗和温度边界测试？
  \item 是否对所有关键错误路径执行过故障注入？
  \item 看门狗、崩溃转储、复位原因和恢复模式是否经过端到端验证？
  \item 测试失败是否能关联到硬件版本、软件版本、工装和原始日志？
  \item 发布配置是否消除了未解释告警，并保留必要的生产诊断能力？
\end{itemize}
